\section{Related Work}
\label{sec:related_work}

\paragraph{Multi-Turn LLM Agents.}
Recent LLM agents extend language models from single-turn text generation to interactive decision making, where the model repeatedly observes an environment, produces an action or tool call, and receives new feedback \citep{yao2022react, yang2024swe, chan2025mle}. These multi-turn agents have been studied in embodied simulators \citep{shridhar2020alfworld}, web navigation \citep{gur2024real, team2025tongyi}, tool use \citep{yao2024tau}, and search-based question answering \citep{jin2025search}. Compared with single-turn tasks, multi-turn agent learning is more challenging because the model must maintain a long interaction history, reason over partially observable states, and recover from its own previous mistakes. An early suboptimal action can change the future observation stream, causing errors to compound across the trajectory \citep{zhang2025agentracer}.

To improve multi-turn agent performance, prior work has mainly explored two training paradigms. The first is multi-turn SFT \citep{teng2024fine, li2025beyond}, which decomposes teacher trajectories into turn-level training examples and trains the model to imitate them. This can improve formatting and local action prediction, but it remains off-policy: the student is optimized on teacher-generated states rather than the histories induced by its own behavior. The second is reinforcement learning with verifiable rewards (RLVR) \citep{shao2024deepseekmath}, where the agent receives reward from the final environment outcome \citep{deepswe2025, wei2025webagent}. While this directly optimizes task success, it typically requires an external verifier or environment-defined reward signal, and the resulting reward is often sparse relative to the long sequence of student-generated turns, making credit assignment difficult~\citep{wang2026openclaw}. These limitations motivate on-policy distillation methods that can provide teacher supervision on student-induced trajectories while avoiding the sparsity of pure outcome-based reward learning.

\paragraph{On-Policy Distillation and Multi-Turn Limitations.}
To mitigate the exposure bias of offline imitation, OPD shifts supervision from teacher-generated trajectories to the student's own generated distribution \citep{agarwal2024policy, lu2025onpolicydistillation}. In OPD, the student first samples trajectories under its current policy, and the teacher then provides token-level feedback on the resulting student-induced states. This aligns training with inference-time behavior and allows the student to learn from states that arise from its own decisions, rather than only from expert demonstrations \citep{zhao2026self, he2026self, hubotter2026reinforcement}.

Despite this advantage, recent analyses have revealed that the token-level logit signal underlying standard OPD is surprisingly brittle. This brittleness is further amplified in multi-turn settings, where local token-level mismatches can compound across turns. First, OPD relies on a narrow token-level overlap between the teacher's and student's plausible next-token distributions \citep{li2026rethinking}. When multiple continuations are semantically valid, dense token-level alignment can over-penalize acceptable alternatives. Second, OPD can reinforce degenerate local patterns when repeated or low-quality prefixes accidentally receive favorable token-level gradients \citep{luo2026demystifying}; in multi-turn agents, this may manifest as repeated thoughts or repeated actions across consecutive turns. Third, multi-turn interaction amplifies distribution shift through compounding errors: student thoughts or actions can push the environment history outside the teacher's effective state distribution, making later teacher supervision less reliable \citep{wang2026tcod, li2025beyond}.

A concurrent line of work, TCOD, addresses this instability by controlling the trajectory depth and gradually expanding the student-controlled horizon through a temporal curriculum \citep{wang2026tcod}. In contrast, \ours{} does not rely on a manually designed curriculum over turn depth. Instead, we delegate the intervention decision to the teacher: for each turn, the teacher determines whether supervision should be skipped, weakly applied, or strongly applied. To further account for compounding errors that may move the student into states where the teacher is less reliable, \ours{} weights the token-level OPD signal by the teacher's predictive confidence. Thus, rather than deciding \emph{how long} to apply OPD, \ours{} decides \emph{where} OPD should be applied and how strongly it should be trusted.